\documentclass{amsart}
\usepackage{amsmath,amssymb,amsthm,hyperref}
\newtheorem{theorem}{Theorem}
\newtheorem{lemma}{Lemma}
\theoremstyle{definition}
\newtheorem{definition}{Definition}
\begin{document}

\title{The infiniteness of the Kleene star of a finite language is in $\mathsf{P}$}
\maketitle

\section{Statement}

A finite language is a finite set $L \subseteq \Sigma^*$ of words over a finite
alphabet $\Sigma$. Its Kleene star is
\[
L^* = \{w_1 w_2 \cdots w_k : k \ge 0,\ w_i \in L\},
\]
where the case $k=0$ yields the empty word $\varepsilon$ (the empty concatenation).
We consider the decision problem

\begin{quote}
INPUT: a finite alphabet $\Sigma$ and a finite list of words
$L = \{w_1,\dots,w_n\} \subseteq \Sigma^*$, each word given explicitly as a string
of symbols; the input size is the total number of symbols used to list the words.

QUESTION: Is $L^*$ infinite?
\end{quote}

We determine its computational complexity.

\begin{theorem}\label{thm:main}
$L^*$ is infinite if and only if $L$ contains at least one nonempty word.
Consequently, the problem ``Is $L^*$ infinite?'' is decidable in linear time
in the input size and therefore lies in $\mathsf{P}$. In fact it is essentially
trivial: it is solvable by a single scan of the input that checks whether any
listed word has length $\ge 1$.
\end{theorem}

\section{Proof}

The key observation is a complete combinatorial characterization of when $L^*$
is infinite. We prove the two directions of the criterion separately.

\begin{lemma}\label{lem:suff}
If $L$ contains a nonempty word, then $L^*$ is infinite.
\end{lemma}

\begin{proof}
Suppose $w \in L$ with $|w| = \ell \ge 1$, where $|\cdot|$ denotes word length.
For every integer $k \ge 0$, the word $w^k = \underbrace{w w \cdots w}_{k}$ is a
concatenation of $k$ words from $L$ (each equal to $w$), hence $w^k \in L^*$ by
the definition of the Kleene star.

Concatenation is additive on lengths: $|uv| = |u| + |v|$ for all words $u,v$.
By induction, $|w^k| = k\ell$. Since $\ell \ge 1$, the map $k \mapsto |w^k| = k\ell$
is strictly increasing in $k$, so the words $\{w^k : k \ge 0\}$ have pairwise
distinct lengths and are therefore pairwise distinct. Thus $L^*$ contains the
infinite set $\{w^k : k \ge 0\}$, and so $L^*$ is infinite.
\end{proof}

\begin{lemma}\label{lem:nec}
If $L$ contains no nonempty word, then $L^*$ is finite; indeed
$L^* \subseteq \{\varepsilon\}$.
\end{lemma}

\begin{proof}
The hypothesis means that every element of $L$ is the empty word $\varepsilon$;
that is, $L \subseteq \{\varepsilon\}$ (so $L$ is either $\emptyset$ or
$\{\varepsilon\}$).

Take any $u \in L^*$. By definition $u = w_1 w_2 \cdots w_k$ with $k \ge 0$ and
each $w_i \in L \subseteq \{\varepsilon\}$, so each $w_i = \varepsilon$. Since
$\varepsilon$ is the identity for concatenation, $u = \varepsilon$ (this also
covers $k = 0$, which gives the empty concatenation $\varepsilon$). Hence every
element of $L^*$ equals $\varepsilon$, i.e. $L^* \subseteq \{\varepsilon\}$.

A subset of the one-element set $\{\varepsilon\}$ is finite; therefore $L^*$ is
finite. (Concretely, $L^* = \{\varepsilon\}$ in both cases $L = \emptyset$ and
$L = \{\varepsilon\}$, since the empty concatenation $k=0$ always produces
$\varepsilon \in L^*$.)
\end{proof}

\begin{proof}[Proof of Theorem~\ref{thm:main}]
Lemmas~\ref{lem:suff} and~\ref{lem:nec} together establish the criterion:
\[
L^* \text{ is infinite} \iff L \text{ contains a word of length } \ge 1.
\]
Indeed, Lemma~\ref{lem:suff} gives the ``$\Leftarrow$'' direction, and
Lemma~\ref{lem:nec} gives the contrapositive of the ``$\Rightarrow$'' direction
(if $L$ has no nonempty word then $L^*$ is finite).

This criterion yields a trivial decision procedure. Read the input, which lists
the words $w_1,\dots,w_n$ as explicit strings of symbols; accept if and only if
some $w_i$ consists of at least one symbol. The empty word is exactly the word
contributing zero symbols to the listing, so detecting a nonempty word amounts
to detecting any symbol attached to some listed word. Scanning the input once
suffices, which costs $O(m)$ time where $m$ is the input size (the total number
of symbols). Hence the problem is solvable in linear time, and in particular it
lies in $\mathsf{P}$.

Finally we address the contrast posed in the problem statement. The problem is
\emph{not} $\mathsf{NP}$-hard unless $\mathsf{P} = \mathsf{NP}$, because it lies
in $\mathsf{P}$: an $\mathsf{NP}$-hardness reduction from any $\mathsf{NP}$-hard
problem would, composed with the polynomial-time decision procedure above, place
that $\mathsf{NP}$-hard problem (and hence all of $\mathsf{NP}$) in $\mathsf{P}$.
Thus, under the standard assumption $\mathsf{P} \neq \mathsf{NP}$, the problem is
not $\mathsf{NP}$-hard. This completes the characterization: the problem is in
$\mathsf{P}$ (indeed solvable in linear time by a single scan).
\end{proof}

\section{Remarks}

\begin{itemize}
\item The result is robust to natural variants of the encoding. Whether $L$ is
given as a set or a list, with or without repetitions, and regardless of how the
alphabet $\Sigma$ is encoded, the only quantity that matters is the existence of
a single listed word of length $\ge 1$. The empty alphabet $\Sigma = \emptyset$
forces $L \subseteq \{\varepsilon\}$, consistent with the criterion.

\item The problem is in fact far below $\mathsf{P}$ in complexity (e.g. it is
$\mathsf{AC}^0$/logspace computable, since it is a disjunction over input
symbols of whether the symbol belongs to some listed word). The theorem only
needs the weaker conclusion that it lies in $\mathsf{P}$, which already settles
the dichotomy raised in the problem.
\end{itemize}

\end{document}
